\documentclass[11pt]{article}
\usepackage[T1]{fontenc}
\usepackage{lmodern}
\usepackage{amsmath,amssymb,amsthm}
\usepackage[margin=1in]{geometry}
\usepackage{microtype}
\usepackage[colorlinks=true,allcolors=blue]{hyperref}
\newtheorem{theorem}{Theorem}
\newtheorem{lemma}[theorem]{Lemma}
\newtheorem{corollary}[theorem]{Corollary}
\newtheorem{proposition}[theorem]{Proposition}
\theoremstyle{remark}\newtheorem{remark}[theorem]{Remark}
\newcommand{\poly}{\operatorname{poly}}
\newcommand{\cc}{\operatorname{cc}}
\newcommand{\ind}[1]{\mathbf 1\{#1\}}
\newcommand{\rhoN}[1]{\frac{#1(\log\log(#1+16))^2}{\sqrt{\log(#1+16)}}}
\title{Connected Dominating Set via Monomial-Weighted SAT Counting\\
\large A linear-variable parity reduction and randomized subexponential search}
\author{}
\date{29 October 2026}
\begin{document}
\maketitle
\begin{abstract}
We construct a width-three CNF with exactly $3n$ variables and
$2n+4m+2$ clauses for a rooted $n$-vertex, $m$-edge graph. For a fixed
selected set $D$ containing the root, its number of extensions is
$2^{u(D)+\cc(G[D])-1}$, where $u(D)$ counts undominated vertices.
Thus its extension parity is precisely the indicator that $D$ is a connected
dominating set. Polynomially bounded isolation weights, attached as formal
monomials, give randomized decision, witness recovery, and optimization.
We prove the required monomial-weight extension of the interpolation route
in the supplied OpenAI manuscript. This yields time
$(n+m+1)^{O(1)}2^{O(n(\log\log n)^2/\sqrt{\log n})}$ and polynomial
space on arbitrary-density graphs, with one-sided error for decision.
The supplied SAT algorithm is an external theorem; the extension here is
not part of the linked Lean scope. Implementations accompanying this note
are exponential small-instance validators, not implementations of that
subexponential algorithm. No historical priority claim is made.
\end{abstract}

\section{Problem, conventions, and outcome}
Let $G=(V,E)$ be a simple undirected graph, $n=|V|$, and $m=|E|$.
A connected dominating set (CDS) is a \emph{nonempty} set $D\subseteq V$
such that $D\cap N[v]\ne\varnothing$ for all $v\in V$ and $G[D]$ is
connected. The decision input also includes an integer $K$, and asks for
$|D|\le K$. A rooted instance additionally requires a specified $r\in D$.
Under our convention the empty graph has no CDS. Inputs with $K<1$ are
negative. Vertices and edges are explicitly represented.
Thresholds are normalized to $0\le K\le n$ before the main algorithm;
reading a longer initial integer representation adds its input length.

\begin{theorem}[Decision and search]\label{thm:main}
Using the interpolation algorithm in the supplied October 25, 2026
manuscript~\cite{oai}, for every positive integer $t$ there is an algorithm
for rooted or unrooted CDS of size at most $K$ with running time
\[
 t(n+m+1)^{O(1)}2^{O(\rhoN{n})}
\]
and $(n+m+\log(t+1)+1)^{O(1)}$ bits of space. It either returns an
independently verifiable CDS satisfying the requirements, or reports
``not detected.'' A negative instance never produces a witness. On a
positive instance, failure to produce a witness has probability at most
$2^{-t}$. The exponential parameter is $n$, even when $m=\Theta(n^2)$.
\end{theorem}

\begin{corollary}[Minimum size]\label{cor:minimum}
Within the same bounds one can return a valid CDS that is minimum with
probability at least $1-2^{-t}$ whenever a CDS exists. Every returned set is
valid; optimality is probabilistic, not independently certified by the
returned set. The rooted version minimizes subject to containing the root.
\end{corollary}

The logarithmic factors in this note deliberately use the interpolation
route. We do not claim that our extension inherits the sharper tensor
contraction bounds of the manuscript. The external theorem and our bridge
are distinct from the elementary CNF identity proved next.

\section{A three-variable-per-vertex parity representation}
Fix $n\ge1$ and $r\in V$. Introduce $x_v,c_v,a_v$ for every vertex $v$.
Here $x$ selects $D$, $c$ marks the right side of a cut of $G[D]$, and
$a_v$ is permitted to be true only when $v$ is undominated. Let $B_{G,r}$
consist of
\begin{align}
 &\neg c_v\vee x_v &&(v\in V),\label{eq:cut-subset}\\
 &x_r,\quad\neg c_r,\label{eq:root}\\
 &\neg c_u\vee\neg x_v\vee c_v,\quad
   \neg c_v\vee\neg x_u\vee c_u &&(\{u,v\}\in E),\label{eq:cut-edge}\\
 &\neg a_v\vee\neg x_w &&(v\in V,\ w\in N[v]).\label{eq:dom-flag}
\end{align}
All variables are declared even if subsequent simplification removes their
occurrences. The construction has exactly $3n$ variables and
\[
 n+2+2m+\sum_v|N[v]|=2n+4m+2
\]
clauses, each of width at most three. Repeating an edge is unnecessary;
the graph representation first deduplicates undirected edges.

Write
\[
 u(D)=|\{v\in V:N[v]\cap D=\varnothing\}|.
\]
\begin{lemma}[Pointwise extension count]\label{lem:fibre}
For every fixed assignment $x=\ind{D}$, the number of satisfying
extensions to $c,a$ is zero if $r\notin D$, and otherwise equals
\[
 2^{u(D)+\cc(G[D])-1}.
\]
Its parity is the indicator that $D$ is a CDS containing $r$.
\end{lemma}
\begin{proof}
The root units force $r\in D$ and place its component on the left.
Equation~\eqref{eq:cut-subset} fixes $c_v=0$ outside $D$.
For an edge inside $D$, the two clauses in~\eqref{eq:cut-edge} force its
endpoints to have the same cut value. An edge with an endpoint outside
$D$ adds no further restriction. Consequently each component other than
the root component chooses its side independently, giving
$2^{\cc(G[D])-1}$ cuts. The clauses in~\eqref{eq:dom-flag} force $a_v=0$
when $v$ is dominated, and otherwise leave $a_v$ free. These independent
choices contribute $2^{u(D)}$. Both exponents are nonnegative, and their
sum vanishes exactly when $u(D)=0$ and $\cc(G[D])=1$.
\end{proof}

\subsection{Generating polynomials and a signed companion}
Set $R_n=2^{\lceil\log_2(2n)\rceil}$, so $2n\le R_n<4n$.
Choose integers $w_v\in\{1,\ldots,R_n\}$ and write
$w(D)=\sum_{v\in D}w_v$. Give the true value of $x_v$ monomial weight
$y z^{w_v}$; all other variable-value weights are one. The resulting
ordinary nonnegative generating polynomial is
\begin{equation}\label{eq:unsigned}
 U_{G,r,w}(y,z)=
 \sum_{(x,c,a)\models B_{G,r}}y^{\sum_vx_v}z^{\sum_vw_vx_v}
 =\sum_{D\ni r}2^{u(D)+\cc(G[D])-1}y^{|D|}z^{w(D)}.
\end{equation}
In particular
\begin{equation}\label{eq:parity}
 [y^kz^W]U_{G,r,w}\bmod2
 =|\{D\ni r:D\text{ is a CDS},\ |D|=k,\ w(D)=W\}|\bmod2.
\end{equation}
The ranges are $0\le k\le n$ and $0\le W\le nR_n<4n^2$, so the entire
coefficient table has polynomial size. Units fixing any subset of the
$x$ variables restrict the outer sum to exactly those choices; the
underlying graph and its domination neighborhoods stay unchanged.

For comparison, assigning weight $-1$ to true $a_v$ gives the signed
polynomial
\begin{equation}\label{eq:signed}
 S_{G,r,w}(y,z)=
 \sum_{\substack{D\ni r\\D\text{ dominates }G}}
 2^{\cc(G[D])-1}y^{|D|}z^{w(D)}.
\end{equation}
Indeed each undominated vertex contributes $1-1=0$ instead of $1+1=2$.
The two polynomials have identical coefficients modulo two.
Equation~\eqref{eq:unsigned} is sufficient for all decision and search
claims here; negative weights and an auxiliary parity chain are unnecessary.
The signed version remains useful as an independent validation identity.

\section{The monomial-weight bridge}\label{sec:bridge}
An ordinary unweighted counting theorem does not by itself give a
generating polynomial at the same Boolean-variable complexity. Encoding
all binary isolation-weight additions by a generic circuit can introduce
$\Theta(n\log n)$ variables. We therefore verify an extension of the
\emph{algorithm}, using its commutative-algebra product interface.

\begin{proposition}[Polynomial-degree monomial weights]\label{prop:weighted}
Fix the clause width $q$ and a constant $d$. Let $F$ be a $q$-CNF with
$N$ declared variables and $M$ clauses. For each variable-value pair choose
\[
 w_i(b)=s_{i,b}y^{\alpha_{i,b}}z^{\beta_{i,b}},\qquad
 s_{i,b}\in\{-1,1\},\quad\alpha_{i,b},\beta_{i,b}\ge0,
\]
where
\[
 A=\sum_i\max_b\alpha_{i,b}\le(N+1)^d,\qquad
 B=\sum_i\max_b\beta_{i,b}\le(N+1)^d.
\]
The interpolation route in~\cite{oai} can be extended to compute
\[
 Z_F(y,z)=\sum_{x\models F}\prod_iw_i(x_i)
\]
exactly in $(N+M+1)^{O_{q,d}(1)}2^{O_{q,d}(\rhoN{N})}$ bit
operations and polynomial space. In particular it computes all
coefficients modulo two within those bounds.
\end{proposition}

\begin{proof}
We track the modifications to the cited algorithm explicitly. Label names
below identify statements in its available TeX source.

\emph{Preserve exact satisfying sets during compression.}
Apply the original disjoint sparsification, with its original parameter
choice of order $(\log\log N)^2/\sqrt{\log N}$, to obtain sparse leaves
on the same declared variables. In each leaf use
\texttt{lem:greedy-hashes} unchanged: all greedy comparisons use
\emph{unweighted} common-zero counts. Every satisfying assignment survives
each check. The final accepted set contains no nonsatisfying assignment,
and is therefore exactly the leaf's satisfying set. Only this final
system receives the monomial factors. Optimizing signed values or
individual coefficients during the greedy phase would not justify this
set equality.

\emph{Enlarge the coefficient algebra, not the Boolean domain.}
Set $\ell=N+2$ and use
\[
 \mathcal A=\mathbb Q[V,y,z]/(V^\ell,y^{A+1},z^{B+1}).
\]
The original $V$ records powers discarded modulo $2^\ell$ in the lifted
equation indicators. The new $y,z$ are external markers. The product
interface \texttt{prop:general-count} accepts a commutative
$\mathbb Q$-algebra. Append one factor of singleton scope $\{i\}$,
\[
 W_i(X_i)=w_i(0)+(w_i(1)-w_i(0))X_i,
\]
per variable. On Boolean inputs this is exactly $w_i(x_i)$.
Give each factor--variable incidence its own input indeterminate, as in
the cited product-degree condition. These factors add at most $N$ to
the total degree in those independent incidence indeterminates, so the
degree remains $O_q(N)$ after truncation in $V$.
Their marker degrees are coefficient data and do not multiply the
interpolation degree in its parameter $T$. Scopes and total incidences
gain only $N$ singleton incidences. The original occurrence branching
and replacement-form construction therefore keep their asymptotic
bounds. A weight of a fixed variable is multiplied outside its branch
sum exactly once; each remaining free variable keeps its unary factor.

\emph{Coefficientwise recovery.}
The original lifted-indicator congruence holds at each Boolean assignment.
Multiplication by its integer monomial weight and summation preserve
that congruence coefficientwise in $y,z$. Thus interpolation, followed by
the original linear map summing retained $V$ coefficients, recovers
$Z_F$ modulo $2^\ell$. Each assignment contributes a single monomial
with coefficient $+1$ or $-1$, so every coefficient has absolute value
at most $2^N$. Since $2^{\ell-1}>2^N$, centered residues recover the
signed integer coefficients uniquely. All actual marker exponents lie
within the stated bounds, so the marker truncations lose no desired
term. Adding the disjoint leaf polynomials gives $Z_F$ without overcounting.

\emph{Bit complexity and space.}
Each coefficient-algebra element has
$\ell(A+1)(B+1)$ rational coordinates, a polynomial number.
At a replacement-form Taylor evaluation, write $d_{\rm den}$ for the
original common denominator of a replaced input and $B_0$ for an upper
bound on the coefficient $1$-norm of its numerator. After multiplication
by $d_{\rm den}$, the additional factor $W_i$ has integer numerator
norm at most $d_{\rm den}+2B_0$. If the original product denominator is
$\Omega$, the enlarged product denominator divides
$\Omega d_{\rm den}^{N}$, uniformly over the assignment summands.
The numerator norm bound gains at most
$(d_{\rm den}+2B_0)^N$. Thus its logarithm increases by only a polynomial
quantity, using the original polynomial bounds on the replacement
forms. Truncated polynomial multiplication preserves the same norm
estimate. Summing exponentially many terms adds only the logarithm of
their number to the bit bound. Apply the original rational interpolation
separately to each marker coordinate; all coordinates share the same
polynomially bounded interpolation data. Consequently each algebra
operation, factor evaluation, and interpolation step has polynomial bit
cost. Sequential leaf and branch traversal retains polynomial space.
The number of Boolean vectors enumerated by the product routine changes
only within the constants of its $O(N)$ degree and polylogarithmic
incidence bounds, giving the asserted exponent.
An empty clause returns the zero polynomial. A formula with no clauses
returns $\prod_i(w_i(0)+w_i(1))$, retaining every declared unused variable.
The finitely many small parameter cases are handled by direct enumeration.
\end{proof}

\subsection{A simpler bridge when only parity is needed}
The CDS algorithm does not need exact integer coefficients. After the
unchanged unweighted greedy phase has produced an exact accepting system,
let $P_1,\ldots,P_h$ be its $\mathbb F_2$ equations, with $h=O(N)$,
and let $L_j$ be their integer lifts of degree at most $q$. For Boolean
$x$,
\[
 1-L_j(x)\equiv\ind{P_j(x)=0}\pmod2.
\]
It therefore suffices to compute
\[
 K(y,z)=\sum_{x\in\{0,1\}^N}
 \prod_{j=1}^h(1-L_j(x))\prod_{i=1}^N W_i(x_i)
\]
and reduce its integer coefficients modulo two. The incidence degree is
at most $qh+N=O_q(N)$. Apply the same product routine over
$\mathbb Q[y,z]/(y^{A+1},z^{B+1})$, with no internal $V$ marker.
The scope argument is unchanged, and the preceding denominator and
norm argument applies with the degree-$q$ factors $1-L_j$ in place of
the lifted indicators. The result at the restored interpolation point
is an integer polynomial, even though intermediate arithmetic is rational.
Thus reduction modulo two takes place only at that point; rational
interpolation is not performed over $\mathbb F_2$.
This gives the same conservative time and space bound with a simpler
final evaluator. Its integer polynomial $K$ need not equal $Z_F$;
only their coefficients modulo two agree. The greedy comparisons
continue to use exact ordinary counts.

This proposition is an extension established in this note using the
supplied algorithm; it is not a claim that the Lean statement~\cite{scope}
already quantifies over such weights. It does not cover arbitrary
exponentially large marker degrees or unrestricted large-coefficient
weights at the same bound.

\section{Isolation, one-sided decision, and recovery}\label{sec:algorithm}
The Isolation Lemma~\cite{isolation,cutcount} states that for a nonempty
family $\mathcal F\subseteq2^V$, independent uniform weights from
$\{1,\ldots,R_n\}$ give a unique minimum-weight member with probability
at least $1-n/R_n\ge1/2$. Each weight uses exactly
$\lceil\log_2(2n)\rceil$ independent random bits, followed by adding one;
there is no unbounded rejection-sampling loop in the stated time bound.

For an unrooted instance take
\[
 \mathcal F_K=\{D\subseteq V:D\text{ is a CDS},\ |D|\le K\}.
\]
One trial samples a single weight vector and evaluates
$U_{G,r,w}$ for every root $r$. It checks for an odd coefficient
$[y^kz^W]U_{G,r,w}$ with $1\le k\le\min(K,n)$.
For a rooted instance evaluate just the specified root and apply
isolation to the family additionally constrained by $r\in D$.

If no feasible CDS exists, equation~\eqref{eq:parity} makes every checked
coefficient even, for every weight vector. If $\mathcal F_K$ is nonempty
and has an isolated minimum-weight member $D_*$, choose any $r\in D_*$.
For $(k_*,W_*)=(|D_*|,w(D_*))$, the only feasible CDS in that coefficient
is $D_*$. Larger-size sets do not contribute to $y^{k_*}$, and the
coefficient is odd. A trial therefore succeeds with probability at least
$1/2$. Independent repetition $t$ times gives the $2^{-t}$ failure bound.
The same weights can be shared by all roots within a trial; independence
between roots is neither used nor required.
The roots must be checked separately: summing their coefficient tables
before reducing modulo two would count each $D$ exactly $|D|$ times,
losing every even-cardinality solution.

\subsection{Recovering a witness without further randomness}
Suppose the coefficient at $(k,W)$ is odd for a root and weight vector.
Maintain the original graph, CNF, weights, and a partial assignment to
$x$. For each next variable $x_v$, compute the same coefficient with an
additional unit $\neg x_v$. If it is odd, keep this unit. Otherwise keep
$x_v$. The parent coefficient is the sum of its two child coefficients,
so exactly one child is odd and this invariant continues. After $n$
choices, a single $D$ remains. Lemma~\ref{lem:fibre} implies that its
odd contribution comes from a CDS containing $r$; the selected
coefficient also ensures $|D|=k$ and $w(D)=W$.

Domination, connectivity, cardinality, and root membership can then be
checked directly in polynomial time. No new isolation experiment is
needed during this self-reduction. In particular, deleting excluded
vertices from the graph would be incorrect: they still require domination.

Proposition~\ref{prop:weighted} applies with $N=3n$, $M=2n+4m+2$,
$A=n$, and $B=nR_n<4n^2$. Root enumeration, coefficient inspection,
repetition, and $n$ recovery calls supply only polynomial factors (and
the explicit factor $t$). This proves Theorem~\ref{thm:main}.

\subsection{Minimum connected dominating set}
Across all roots and trials, keep an odd coefficient with smallest
cardinality index and recover its witness at the end. Let $k_*$ be the
true optimum and apply the Isolation Lemma to the nonempty family of
all CDSs of size exactly $k_*$. With probability at least $1/2$ per
trial, some root has an odd coefficient at cardinality $k_*$. No
coefficient at a smaller cardinality can be odd. Consequently the best
recorded cardinality equals $k_*$ with probability at least $1-2^{-t}$.
This proves Corollary~\ref{cor:minimum}, including its rooted variant.

\section{What transfers beyond this example}
The reusable interface is a fixed-width CNF with $O(n)$ declared
variables and polynomially many clauses, whose variables include the
solution's characteristic vector $x$. If, for each fixed $x$, its
auxiliary-extension parity equals the desired solution predicate, then
the monomial-weight bridge and isolation give the same kind of
randomized subexponential decision and recovery algorithm. Unique
extensions are unnecessary; even multiplicities are useful.
CDS supplies such an interface even on dense graphs. Extending this
principle to another problem still requires an explicit representation
with the promised number of Boolean variables. It does not automatically
give analogous bounds for Hamiltonicity or feedback vertex set.

\section{Prior work, validation, and limitations}
The consistent-cut multiplicity and isolation step are established
Cut\&Count ingredients~\cite{cutcount}; this note does not claim them
as new. Our construction represents the domination filter with one
free auxiliary bit per undominated vertex and performs the required
coefficient computations through the supplied SAT algorithm.
The linear-variable accounting and the monomial-weight bridge are the
specific items proved and audited here. Historical priority remains
unconfirmed.

The known linear-measure existential monadic second-order framework
contains several local graph problems~\cite{linear}, but connectivity
is not itself definable by existential quantification over only unary
relations on the original graph signature~\cite{monadic}. In particular,
a CDS verifier in that language, with a free unary predicate $D$, would
define graph connectivity after setting $D=V$. This observation concerns
that unaugmented signature; it does not exclude other preprocessing,
encodings, or algorithms.

The accompanying \texttt{code/connected\_domination.py} implements the
CNF and a small exact generating-polynomial DPLL, with signed and
unsigned modes. The independent verifier compares its coefficients
against direct graph definitions, tests conditioning on selection bits,
and checks recovered witnesses. Run
\texttt{code/verify\_connected\_domination.py} to generate a JSON report
and text summary under \texttt{workspace/}. These experiments check the representation and the
implemented finite algorithm; they do not execute or benchmark the
external subexponential counting algorithm.

Neither polynomial above directly counts CDSs over the integers:
disconnected selected sets have positive even multiplicity. Taking
coefficients modulo two yields parity counts, and isolation converts
these into randomized decision and search. Deterministic existence,
exact CDS counting, and a practical implementation of the extended
OpenAI counter remain separate tasks.

\begin{thebibliography}{9}
\bibitem{oai} OpenAI.
\emph{The Exponential Time Hypothesis is False}. October 25, 2026 revision.
\href{https://raw.githubusercontent.com/openai/math/refs/heads/main/preprints/The-Exponential-Time-Hypothesis-is-False-October-25-2026/build/main_full_version.tex}{Full manuscript source}.
\bibitem{scope} OpenAI.
\href{https://github.com/openai/math/blob/main/lean/docs/378.md}{Lean scope document 378}.
\bibitem{cutcount} M. Cygan, J. Nederlof, M. Pilipczuk, M. Pilipczuk,
J. M. M. van Rooij, J. O. Wojtaszczyk.
\emph{Solving connectivity problems parameterized by treewidth in single exponential time}.
\href{https://arxiv.org/abs/1103.0534}{arXiv:1103.0534}, 2011.
\bibitem{isolation} K. Mulmuley, U. V. Vazirani, V. V. Vazirani.
\emph{Matching is as easy as matrix inversion}.
Combinatorica 7, 105--113, 1987.
\bibitem{linear} R. Ganian, R. de Haan, I. Kanj, S. Szeider.
\emph{On Existential MSO and its Relation to ETH}.
\href{https://drops.dagstuhl.de/entities/document/10.4230/LIPIcs.MFCS.2016.42}{MFCS 2016, 42:1--42:15}.
\bibitem{monadic} R. Fagin, L. J. Stockmeyer, M. Y. Vardi.
\emph{On Monadic NP vs. Monadic co-NP}.
Information and Computation 120, 78--92, 1995.
\href{https://s3.us.cloud-object-storage.appdomain.cloud/res-files/500-infcom95.pdf}{Author-hosted article}, Theorem 5.1.
\end{thebibliography}
\end{document}
